\documentclass[UTF8,a4paper,11pt]{ctexart}

\usepackage{amsmath,amssymb,bm}
\usepackage{booktabs}
\usepackage{geometry}
\usepackage{array}
\usepackage{float}
\usepackage{graphicx}
\usepackage{caption}
\usepackage{hyperref}
\usepackage{enumitem}

\geometry{left=2.25cm,right=2.25cm,top=2.15cm,bottom=2.15cm}
\setlength{\parindent}{2em}
\linespread{1.10}
\captionsetup{font=small}
\setlist{nosep}
\hypersetup{colorlinks=true,linkcolor=black,citecolor=black,urlcolor=blue}

\begin{document}

\section{问题一：预热阶段药材内部温度与水分浓度分布}

\subsection{问题分析与模型思路}

热风干燥属于典型的非稳态热质传递过程。外界热空气首先通过药材表面进行对流换热，
热量再沿温度梯度向内部传递；与此同时，表面水分向外界迁移后形成内部水分梯度，
推动水分由中心区域向表面扩散。中药材干燥研究通常将内部水分迁移和表面蒸发概括为
“导湿—给湿”两个基本环节\cite{wang2024}。

本题要求计算多个时刻、多个径向位置处的温度和干基含水率，
因此仅采用平均温度或平均含水率的集中参数模型无法满足要求。
Da Silva 等针对圆柱状物料建立了一维圆柱坐标热质传递模型，
并采用第三类边界条件及全隐式有限体积法求解\cite{dasilva2014}。
结合本题药材几何形状、输出位置和题设参数，本文建立一维径向非稳态传热—传质模型。

问题一给出的 $\rho,c_p,k$ 为常数，而水分扩散系数仅依赖局部含水率，
因此本问作为全文基础模型，采用固定几何、常热物性和弱耦合假设，
分别求解温度场与水分场；后续问题二再解除常物性假设并引入状态相关耦合。

\subsection{模型假设与一维化依据}

药材长度 $L=0.25~\mathrm m$，半径 $R=0.02~\mathrm m$，有
\[
\frac{L}{R}=12.5,
\qquad
\frac{\text{两端总面积}}{\text{侧面积}}
=\frac{R}{L}=0.08.
\]
因此对于圆柱主体中段，侧表面交换占主导，本文忽略端部附近的轴向梯度，
仅研究径向状态
\[
T=T(r,t),\qquad C=C(r,t),\qquad 0\le r\le R.
\]

进一步计算初始 Biot 数：
\[
Bi_T=\frac{h_TR}{k}\approx1.389,
\qquad
D(C_0)\approx4.94\times10^{-9}\ \mathrm{m^2/s},
\qquad
Bi_C=\frac{h_mR}{D(C_0)}\approx3.24.
\]
二者均明显大于集中参数模型常用的 $0.1$ 判据，
说明药材内部径向温度梯度和水分梯度不可忽略，采用空间分布模型是必要的。

在此基础上作如下假设：

\begin{enumerate}
    \item 药材视为均质、各向同性连续介质，初始温度和干基含水率均匀；
    \item 前 $1800~\mathrm s$ 内圆柱尺寸保持不变；
    \item 问题一中 $\rho,c_p,k$ 取附录2给定常数；
    \item 附件1中的烘房水分浓度作为表面传质的等效外界浓度驱动力；
    \item 暂不考虑蒸发潜热、辐射换热及内部热源对温度场的反馈。
\end{enumerate}

其中，第4条属于题设参数体系下的宏观等效边界；
第5条用于建立基础模型，并不意味着潜热在真实干燥过程中一定可以忽略。
已有圆柱物料干燥研究表明，潜热在部分材料的温度预测中可能具有明显影响
\cite{dasilva2014}。

\subsection{物理参数与数据来源}

\begin{table}[H]
\centering
\caption{问题一主要参数与来源}
\begin{tabular}{clcll}
\toprule
符号 & 物理意义 & 数值或关系式 & 单位 & 来源\\
\midrule
$L$ & 药材长度 & $0.25$ & m & 题目\\
$R$ & 药材半径 & $0.02$ & m & 题目\\
$T_0$ & 初始温度 & $28$ & $^\circ$C & 题目\\
$C_0$ & 初始干基含水率 & $2.55$ & kg/kg & 题目\\
$\rho$ & 密度 & $820$ & kg/m$^3$ & 附录2\\
$c_p$ & 比热容 & $2600$ & J/(kg$\cdot$K) & 附录2\\
$k$ & 导热系数 & $0.36$ & W/(m$\cdot$K) & 附录2\\
$h_T$ & 对流换热系数 & $25$ & W/(m$^2\cdot$K) & 附录2\\
$h_m$ & 等效传质系数 & $8\times10^{-7}$ & m/s & 附录2\\
$D(C)$ & 有效水分扩散系数 & $7\times10^{-9}e^{-0.89/C}$ & m$^2$/s & 附录2\\
\bottomrule
\end{tabular}
\end{table}

附件1给出的烘房温度 $T_\infty(t)$ 和水分浓度 $C_\infty(t)$
作为时变边界输入。由于采样间隔为 $60~\mathrm s$，
而数值积分步长远小于该间隔，计算中采用分段线性插值：
\[
B(t)=B(\tau_j)+
\frac{t-\tau_j}{\tau_{j+1}-\tau_j}
\left[B(\tau_{j+1})-B(\tau_j)\right],
\]
其中 $B$ 表示 $T_\infty$ 或 $C_\infty$。
需要区分：$60~\mathrm s$ 为数据采样间隔，$1~\mathrm s$ 为结果输出间隔，
而数值积分步长 $\Delta t$ 由离散误差检验确定。

\subsection{控制方程与边界条件}

\subsubsection{温度场}

根据能量守恒和 Fourier 导热定律，
\begin{equation}
\boxed{
\rho c_p\frac{\partial T}{\partial t}
=
\frac1r
\frac{\partial}{\partial r}
\left(
rk\frac{\partial T}{\partial r}
\right)
},
\qquad 0<r<R.
\label{eq:heat}
\end{equation}

热扩散率
\[
\alpha=\frac{k}{\rho c_p}
\approx1.6886\times10^{-7}\ \mathrm{m^2/s}.
\]

\subsubsection{水分场}

以干基含水率
\[
C=\frac{m_w}{m_d}
\]
为状态变量，采用有效扩散模型描述内部复杂水分迁移过程：
\begin{equation}
\boxed{
\frac{\partial C}{\partial t}
=
\frac1r
\frac{\partial}{\partial r}
\left[
rD(C)\frac{\partial C}{\partial r}
\right]
},
\qquad 0<r<R.
\label{eq:mass}
\end{equation}

其中
\[
D(C)=7\times10^{-9}\exp\left(-\frac{0.89}{C}\right)
\ \mathrm{m^2/s}.
\]
由于 $D$ 随局部含水率变化，必须保留其通量散度形式，
不能将 $D$ 简单移至空间微分算子之外。

\subsubsection{初始条件与第三类边界}

初始状态为
\[
T(r,0)=28^\circ\mathrm C,
\qquad
C(r,0)=2.55.
\]

轴线处满足对称条件
\[
\left.\frac{\partial T}{\partial r}\right|_{r=0}=0,
\qquad
\left.\frac{\partial C}{\partial r}\right|_{r=0}=0.
\]

侧表面采用第三类边界：
\begin{equation}
-k\left.\frac{\partial T}{\partial r}\right|_{r=R}
=
h_T\left(T_s-T_\infty\right),
\label{eq:heatbc}
\end{equation}
\begin{equation}
-D(C_s)\left.\frac{\partial C}{\partial r}\right|_{r=R}
=
h_m\left(C_s-C_\infty\right),
\label{eq:massbc}
\end{equation}
其中
\[
T_s=T(R,t),\qquad C_s=C(R,t).
\]

需要指出，附件1中的烘房水分浓度与药材干基含水率并非严格相同的热力学量，
本文依据题设给定的 $h_m$ 将其作为等效外界浓度驱动力使用，
因此式~\eqref{eq:massbc} 属于宏观经验型 Robin 边界。

\subsection{有限体积离散与数值求解}

式~\eqref{eq:heat}--\eqref{eq:mass} 包含圆柱几何项、时变 Robin 边界，
且水分扩散系数具有非线性。
本文采用径向有限体积法，并使用后向 Euler 全隐式时间推进。
有限体积法直接对控制体积分，可保持“储存量变化=界面净通量”的守恒结构，
同时便于处理圆柱几何与第三类边界\cite{dasilva2014}。

将 $0\le r\le R$ 均匀划分为 $N$ 个环形控制体，
第 $i$ 个控制体的体积和界面面积分别为
\[
V_i=\pi L\left(r_{i+1/2}^2-r_{i-1/2}^2\right),
\qquad
A_{i\pm1/2}=2\pi Lr_{i\pm1/2}.
\]

\IfFileExists{问题1_有限体积离散示意图_论文风格.png}{
\begin{figure}[H]
\centering
\includegraphics[width=0.78\textwidth]{问题1_有限体积离散示意图_论文风格.png}
\caption{圆柱药材径向有限体积离散示意图}
\end{figure}
}{
\begin{figure}[H]
\centering
\fbox{\parbox[c][3.3cm][c]{0.78\textwidth}{\centering
此处插入“问题1\_有限体积离散示意图\_论文风格.png”}}
\caption{圆柱药材径向有限体积离散示意图}
\end{figure}
}

定义
\[
q=-k\frac{\partial T}{\partial r},
\qquad
f=-D(C)\frac{\partial C}{\partial r}.
\]
对控制体积分得到
\[
\rho c_pV_i
\frac{T_i^{n+1}-T_i^n}{\Delta t}
=
A_{i-1/2}q_{i-1/2}^{n+1}
-
A_{i+1/2}q_{i+1/2}^{n+1},
\]
\[
V_i
\frac{C_i^{n+1}-C_i^n}{\Delta t}
=
A_{i-1/2}f_{i-1/2}^{n+1}
-
A_{i+1/2}f_{i+1/2}^{n+1}.
\]

内部界面的水分扩散系数采用调和平均
\[
D_{i+1/2}
=
\frac{2D(C_i)D(C_{i+1})}
{D(C_i)+D(C_{i+1})},
\]
相当于将界面两侧半控制体的扩散阻力串联。

对于最外层控制体中心 $P$，节点与真实表面相距 $\Delta r/2$。
因此将半单元内部传递阻力与外部对流阻力串联：
\[
q_R=
\frac{T_P-T_\infty}
{\dfrac{\Delta r}{2k}+\dfrac1{h_T}},
\qquad
T_s=T_\infty+\frac{q_R}{h_T},
\]
\[
f_R=
\frac{C_P-C_\infty}
{\dfrac{\Delta r}{2D_s}+\dfrac1{h_m}},
\qquad
C_s=C_\infty+\frac{f_R}{h_m}.
\]

温度方程形成线性三对角系统，可由 Thomas 追赶法求解；
水分方程因 $D=D(C)$ 非线性，在每个时间层采用 Picard 迭代，
收敛判据取
\[
\max\left(
\|C^{(m+1)}-C^{(m)}\|_\infty,
|C_s^{(m+1)}-C_s^{(m)}|
\right)<10^{-10}.
\]

输出时，轴心值采用对称条件下的二阶重构
\[
u(0)\approx\frac{9u_1-u_2}{8},
\]
内部指定半径采用线性插值，表面值由 Robin 边界关系反算。
该处理避免将最外层单元中心误当作真实表面。

\subsection{数值可靠性检验}

最终计算采用
\[
N=1600,
\qquad
\Delta r=1.25\times10^{-5}\ \mathrm m,
\qquad
\Delta t=0.015625\ \mathrm s,
\]
Picard 收敛容限为 $10^{-10}$，
最终计算中每个时间层最多4次迭代即可收敛。

以题目规定的全部输出位置和时刻为比较集合，
定义最大绝对差
\[
E_T=\max|T^{(a)}-T^{(b)}|,
\qquad
E_C=\max|C^{(a)}-C^{(b)}|.
\]

\begin{table}[H]
\centering
\caption{问题一数值可靠性检验}
\begin{tabular}{lccc}
\toprule
检验项目 & 对比设置 & $E_T$/($^\circ$C) & $E_C$/(kg/kg)\\
\midrule
空间网格细化 &
$N=800\rightarrow1600$ &
$1.02\times10^{-6}$ &
$1.522\times10^{-5}$\\
时间步细化 &
$0.03125\rightarrow0.015625$ s &
$2.986\times10^{-5}$ &
$6.39\times10^{-6}$\\
Picard 容差收紧 &
$10^{-8}\rightarrow10^{-10}$ &
$0$ &
$2.15\times10^{-13}$\\
\bottomrule
\end{tabular}
\end{table}

此外，最大相对离散平衡残差为
\[
\varepsilon_T^{\mathrm{bal}}
=3.919\times10^{-7},
\qquad
\varepsilon_C^{\mathrm{bal}}
=3.103\times10^{-11}.
\]
空间离散、时间离散和非线性迭代误差均已显著小于最终结果显示精度对应的量级。
上述检验说明数值求解稳定且离散方程闭合，
但不能替代真实内部温度或含水率实验数据对物理模型的验证。

\subsection{结果与物理解释}

\begin{table}[H]
\centering
\caption{30 min 内药材径向温度分布（$^\circ$C）}
\begin{tabular}{rrrrrr}
\toprule
时间/s & 0 cm & 0.5 cm & 1.0 cm & 1.5 cm & 2.0 cm\\
\midrule
100  & 28.0001 & 28.0003 & 28.0040 & 28.0327 & 28.1801\\
300  & 28.0408 & 28.0635 & 28.1514 & 28.3681 & 28.8489\\
600  & 28.4534 & 28.5361 & 28.8040 & 29.3159 & 30.1652\\
900  & 29.3243 & 29.4583 & 29.8755 & 30.6162 & 31.7304\\
1200 & 30.5427 & 30.7098 & 31.2224 & 32.1126 & 33.4277\\
1500 & 31.9957 & 32.1868 & 32.7660 & 33.7463 & 35.1203\\
1800 & 33.5754 & 33.7720 & 34.3642 & 35.3621 & 36.7856\\
\bottomrule
\end{tabular}
\end{table}

\begin{table}[H]
\centering
\caption{30 min 内药材径向水分浓度分布（kg/kg，干基）}
\begin{tabular}{rrrrrr}
\toprule
时间/s & 0 cm & 0.5 cm & 1.0 cm & 1.5 cm & 2.0 cm\\
\midrule
100  & 2.5500 & 2.5500 & 2.5500 & 2.5500 & 2.2470\\
300  & 2.5500 & 2.5500 & 2.5500 & 2.5492 & 2.0508\\
600  & 2.5500 & 2.5500 & 2.5500 & 2.5353 & 1.8770\\
900  & 2.5500 & 2.5500 & 2.5497 & 2.5045 & 1.7547\\
1200 & 2.5500 & 2.5500 & 2.5482 & 2.4646 & 1.6586\\
1500 & 2.5500 & 2.5499 & 2.5445 & 2.4206 & 1.5787\\
1800 & 2.5500 & 2.5497 & 2.5383 & 2.3755 & 1.5102\\
\bottomrule
\end{tabular}
\end{table}

\IfFileExists{问题1_径向温度分布_重算.png}{
\begin{figure}[H]
\centering
\includegraphics[width=0.78\textwidth]{问题1_径向温度分布_重算.png}
\caption{不同预热时刻的径向温度分布}
\end{figure}
}{}

\IfFileExists{问题1_径向水分浓度分布_重算.png}{
\begin{figure}[H]
\centering
\includegraphics[width=0.78\textwidth]{问题1_径向水分浓度分布_重算.png}
\caption{不同预热时刻的径向水分浓度分布}
\end{figure}
}{}

结果表明，温度变化已深入圆柱主体区域，而水分变化仍主要集中在近表层。
这一差异可以由特征尺度解释。

在 $t=1800~\mathrm s$ 时，
\[
\ell_T\sim\sqrt{\alpha t}
\approx1.74~\mathrm{cm},
\]
已接近半径 $R=2~\mathrm{cm}$；而
\[
\ell_C\sim\sqrt{D(C_0)t}
\approx0.30~\mathrm{cm},
\]
仅约占半径的15\%。因此
\[
\ell_T\gg \ell_C.
\]
这与数值结果中“中心温度已明显升高，而中心含水率仍接近初值”的现象一致，
说明预热阶段具有显著的热—湿时间尺度分离。

\IfFileExists{问题1_温度时空云图.png}{
\begin{figure}[H]
\centering
\includegraphics[width=0.76\textwidth]{问题1_温度时空云图.png}
\caption{问题一温度场时空分布}
\end{figure}
}{}

\IfFileExists{问题1_水分浓度时空云图.png}{
\begin{figure}[H]
\centering
\includegraphics[width=0.76\textwidth]{问题1_水分浓度时空云图.png}
\caption{问题一水分浓度场时空分布}
\end{figure}
}{}

温度云图中等温线由表层逐步向轴心推进，
而水分云图中低含水区在30 min 内仍主要停留在近表面区域，
进一步说明热量传递快于水分扩散。

\subsection{模型讨论与问题二衔接}

问题一模型的主要限制来自三方面：
其一，一维近似主要适用于圆柱中段，端面附近仍可能存在二维效应；
其二，附件1中的烘房水分浓度采用等效 Robin 边界解释；
其三，基础模型暂未考虑蒸发潜热反馈。
此外，几何收缩将在问题四中利用题目给定的尺寸变化进一步处理。

问题一由此建立了固定几何、常物性的径向热质传递基础框架。
完整干燥过程中，含水率持续下降会导致
$\rho$、$c_p$、$k$ 以及 $D$ 随状态变化。
因此问题二将在同一数值框架上解除常物性假设，
利用附录3给出的状态相关关系建立变物性热质耦合模型：
\[
\boxed{
M_1:\ \text{固定几何、常物性基础模型}
}
\quad\Longrightarrow\quad
\boxed{
M_2:\ \text{固定几何、变物性热质耦合模型}
}.
\]

\begin{thebibliography}{9}

\bibitem{wang2024}
王乐意, 李长河, 刘明政, 等.
中药材干燥技术与装备研究现状[J].
农业工程学报, 2024, 40(2): 1--28.
DOI: 10.11975/j.issn.1002-6819.202306104.

\bibitem{dasilva2014}
DA SILVA W P, E SILVA C M D P S, GAMA F J A.
Estimation of thermo-physical properties of products with cylindrical shape during drying:
The coupling between mass and heat[J].
Journal of Food Engineering, 2014, 141: 65--73.
DOI: 10.1016/j.jfoodeng.2014.05.010.

\end{thebibliography}

\end{document}
